\section{Builtins}
\label{sec:builtins}

The following functions are available without declaration. ``Trap'' marks
a runtime error (Section~\ref{sec:errors}).

\subsection{Sequences and text}

\begin{description}[style=nextline,leftmargin=5.2em]
\item[\ic{len(x)}] Length of a string (bytes), array (elements) or map
  (entries) as \ic{long}. Records are rejected (a record's field count
  is part of its type).
\item[\ic{range(n)}] \ic{[0, 1, ..., n-1]}.
\item[\ic{trim(s)}] Strip leading/trailing space, tab, CR, LF.
\item[\ic{join(xs, sep)}] Concatenate \ic{[str]} with a separator.
\item[\ic{split(s, sep)}] Split on every occurrence of a non-empty
  separator; returns \ic{[str]} (empty for empty separator).
\item[\ic{index\_of(s, sub)}] First index of \ic{sub} in \ic{s}, or
  \ic{-1}.
\item[\ic{last\_index\_of(s, sub)}] Last index, or \ic{-1}.
\item[\ic{ord(s)}] First byte as \ic{long}, or \ic{-1} for empty.
\item[\ic{chr(n)}] One-byte string of \ic{n} (masked to 8 bits).
\item[\ic{byte\_at(s, i)}] Byte at \ic{i} as \ic{long}; traps out of
  bounds. Cheap per-byte access (no copy).
\end{description}

\noindent There is no \ic{substr}: a substring is a slice,
\ic{s[i:i + n]} (Section~\ref{sec:slices}).

\subsection{String methods}
\label{sec:string-methods}

A \ic{str} value takes methods: \ic{"str 1 2 3".split(" ")} is
\ic{["str", "1", "2", "3"]}. A method call is checked like a function
call (argument count and types). Method names are not reserved: a
function, a variable or a record field may share one, and \ic{r.f(x)} on
a record still calls the field. \ic{s.length} is a property, so
\ic{s.length()} is an error. Widths and positions count bytes.

\begin{description}[style=nextline,leftmargin=5.2em]
\item[Builtins as methods] \ic{s.split(sep)}, \ic{s.trim()} (also
  \ic{s.strip()}), \ic{s.index\_of(sub)}, \ic{s.last\_index\_of(sub)},
  \ic{s.ord()}, \ic{s.byte\_at(i)}, \ic{s.tonum()} and \ic{s.bytes()} are
  the builtins of the same name with \ic{s} as the first argument.
\item[\ic{sep.join(xs)}] \ic{join(xs, sep)}: the receiver is the
  separator.
\item[\ic{s.trim\_start()}, \ic{s.trim\_end()}] Strip space, tab, CR and
  LF at one end only (also spelled \ic{lstrip}, \ic{rstrip}).
\item[\ic{s.pad\_left(w, fill)}, \ic{s.pad\_right(w, fill)}, \ic{s.center(w, fill)}]
  Widen \ic{s} to \ic{w} bytes with \ic{fill} (default \ic{" "}, repeated
  and cut to fit); a longer \ic{s} is returned unchanged. \ic{center}
  puts the odd byte on the right.
\item[\ic{s.starts\_with(p)}, \ic{s.ends\_with(p)}, \ic{s.contains(p)}]
  \ic{bool}; \ic{""} is a prefix, suffix and part of every string.
\item[\ic{s.count(p)}] Non-overlapping occurrences, left to right
  (\ic{0} for an empty \ic{p}).
\item[\ic{s.replace(p, r)}] Every non-overlapping occurrence of \ic{p},
  left to right, replaced by \ic{r}; an empty \ic{p} replaces nothing.
\item[\ic{s.upper()}, \ic{s.lower()}] ASCII case mapping; other bytes are
  kept.
\item[\ic{s.repeat(n)}] \ic{s} written \ic{n} times; \ic{""} when
  \ic{n <= 0}.
\item[\ic{s.lines()}] Split at each LF, dropping a CR just before it; a
  final newline does not start an empty line, so \ic{"a\bs nb\bs n".lines()}
  is \ic{["a", "b"]} and \ic{"".lines()} is \ic{[]}.
\end{description}

\subsection{Conversions}
\label{sec:builtin-tostring}

\begin{description}[style=nextline,leftmargin=5.2em]
\item[\ic{tostring(v)}] String form: numbers as decimal (floats with six
  decimals), \ic{true}/\ic{false}, strings unchanged, options as their
  payload or \ic{"none"}, arrays as \ic{"[e0, e1, ...]"}, records as
  \ic{"\{k: v, ...\}"}, maps as \ic{"\{"k": v, ...\}"} (string keys
  quoted).
\item[\ic{tonum(s)}] Parse a decimal integer (\ic{0} on garbage).
\item[\ic{typeof(v)}] The static type as a printable value
  (\ic{typeof(none)} is \ic{"none"}).
\end{description}

\subsection{Binary data}
\label{sec:builtin-bytes}

\begin{description}[style=nextline,leftmargin=5.2em]
\item[\ic{bytes(s)}] \ic{[byte]} of the string's raw bytes.
\item[\ic{from\_bytes(b)}] String from a \ic{[byte]} or fixed \ic{[byte; N]} array.
\end{description}

\subsection{Options}
\label{sec:builtin-option-b}

\begin{description}[style=nextline,leftmargin=5.2em]
\item[\ic{some(v)}] A present option holding \ic{v}.
\item[\ic{is\_some(v)}] \ic{true} unless \ic{v} is \ic{none}.
\item[\ic{unwrap(v)}] The payload; \textbf{traps} on \ic{none}. The
  postfix form \ic{v!} is the same operation.
\end{description}

\subsection{Input/output and the environment}

\begin{description}[style=nextline,leftmargin=5.2em]
\item[\ic{print(x)}] Print a value followed by a newline to stdout.
\item[\ic{puts(s)}] Write a string to stdout, no newline added.
\item[\ic{eputs(s)}] Same, to stderr (keeps \ic{--time} output and
  program output apart).
\item[\ic{readline()}] One line from stdin without its newline, or
  \ic{none} at end of input (\ic{str?}). Stdin is buffered (16\,KB per
  \ic{read(2)}); before \ic{exec()} starts a child, the unread part is
  given back when stdin is a file, so the child reads on from the next
  line. On a pipe, lines already buffered stay with the program, as with
  C stdio.
\item[\ic{clear\_eof()}] No-op: \ic{readline()} reads again whenever
  its buffer is empty, so end of input is never latched. Present so the
  same source runs on every compiler.
\item[\ic{read\_file(path)}] File contents as \ic{str?}: \ic{none} when
  the file cannot be read.
\item[\ic{write\_file(path, s)}] Write a string; \ic{bool} success.
\item[\ic{write\_bytes(path, b, mode)}] Write \ic{[byte]} binary-safe
  and chmod (\ic{mode == -1} skips chmod --- any other value is passed
  to \ic{fchmodat} verbatim, so \ic{493} is \ic{0755}); \ic{bool} success.
\item[\ic{argv()}] \ic{[str]} of the program's arguments.
\item[\ic{exit(code)}] Terminate with an exit code.
\item[\ic{clock\_us()}] Monotonic microseconds as \ic{long}.
\item[\ic{heap\_used()}] Bytes allocated so far.
\item[\ic{sleep(ms)}] Suspend the process (\ic{void}).
\item[\ic{syscall(...)}] Raw Linux syscall wrapper; the first argument
  is the syscall number.
\item[\ic{exec(cmd)}] Run a shell command, return its exit code.
  Removed from the binary entirely by \ic{--no-exec} (then calling it
  is a compile error).
\end{description}

\subsection{Pointers}
\label{sec:builtin-map}

\begin{description}[style=nextline,leftmargin=5.2em]
\item[\ic{ptr\_addr(p)}] The address of a pointer as \ic{long}.
\end{description}

\noindent Mutating and keying are symbolic, not functional: arrays use
\ic{+=} / \ic{*=}, records use \ic{r.f}, maps use \ic{m['k']} and
\ic{m['k'] = v}. The legacy functions \ic{push}, \ic{prepend},
\ic{map\_get} and \ic{map\_set} do not exist.
